\documentclass[12pt,a4paper]{article}
\usepackage{ugmath}
\usepackage{float}
\usepackage{placeins}
\usepackage{array}
\usepackage[labelfont=bf,labelsep=space]{caption}
\newcommand{\paren}[1]{\left( #1 \right)}
\newcommand{\alumno}{Ricardo León Martínez}
\newcommand{\materia}{Elementos de Probabilidad y Estadistica}
\newcommand{\profesor}{Claudia Reynoso Alcántara}
\newcommand{\tarea}{Tarea 5}
\newcommand{\fecha}{25/3/2026}

\begin{document}

    \begin{center}
        {\large\textbf{UNIVERSIDAD DE GUANAJUATO}}\\[0.3cm]
        {\normalsize\textbf{DIVISIÓN DE CIENCIAS NATURALES Y EXACTAS}}\\
        {\normalsize\textbf{CAMPUS GUANAJUATO}}\\[1cm]

        {\Large\textbf{\tarea\ (\materia)}}\\[1cm]
    \end{center}

    \noindent
    \textbf{Nombre:} \alumno 
    \hfill 
    \textbf{Fecha:} \fecha 
    \hfill 
    \textbf{Calificación:} \rule{3cm}{0.4pt}

    \vspace{0.3cm}

    \begin{mdframed}[style=mdbluebox,frametitle={Ejercicio 1}]
        Sea $\{A_k\}_{k=1}^\infty$ una sucesión de eventos en un espacio medible $(\Omega,\mathcal{F})$. Para cada entero $n \geq 1$, consideramos los conjuntos
        \[
        B_n = \inf_{k \geq n} A_k := \bigcap_{k=n}^{\infty} A_k
        \quad \text{y} \quad
        C_n = \sup_{k \geq n} A_k := \bigcup_{k=n}^{\infty} A_k .
        \]
        Notemos que $\{B_n\}_{n=1}^\infty$ es creciente mientras que $\{C_n\}_{n=1}^\infty$ es decreciente, esto es, 
        $B_1 \subset B_2 \subset B_3 \subset \dots$ mientras que $C_1 \supset C_2 \supset C_3 \supset \dots$ . Podemos también considerar
        \[
        B = \lim_{n \to \infty} B_n := \bigcup_{n=1}^{\infty} B_n 
        \quad \text{y} \quad
        C = \lim_{n \to \infty} C_n := \bigcap_{n=1}^{\infty} C_n .
        \]
        Los conjuntos $B$ y $C$ se conocen como el límite inferior y el límite superior de la sucesión $\{A_k\}_{k=1}^\infty$, y se denotan por $\liminf_k A_k$ y $\limsup_k A_k$, respectivamente. Demuestre que
        \begin{enumerate}
            \item[(a)] $B_{n},C_{n},B,C$ son eventos, es decir, que pertenecen a la
            $\sigma$-algebra $\mathcal{F}$.
            \item[(b)] $B=\left\{\omega\in\Omega:\begin{matrix}
            \omega\in A_{k}\text{ para todos los valores de $k$ salvo, posiblemente}\\
            \text{por un número finito de los mismos}
            \end{matrix}\right\}$
            \item[(c)] $C=\{\omega\in\Omega:\omega\in A_{k}\text{ para un número infinito de valores de $k$}\}$
        \end{enumerate}
    \end{mdframed}

    \begin{proof}
        \begin{enumerate}
            \item[(a)] $B_{n},C_{n},B,C$ son eventos, es decir, que pertenecen a la
            $\sigma$-algebra $\mathcal{F}$.\\
            Dado que $\mathcal{F}$ es una $\sigma$-algebra, es cerrada bajo uniones e intersecciones
            numebrales. Para cada $n$ $B_{n}=\bigcap_{k=n}^{\infty}A_{k}$ es una interseccion numerable de
            elementos de $\mathcal{F}$, por lo tanto $B_{n}\in\mathcal{F}$. Analogamente, $C_{n}=\bigcup_{k=n}^{\infty}A_{k}$
            es una union numerable, luego $C_{n}\in\mathcal{F}$. Finalmente, $B=\bigcup_{n=1}^{\infty}B_{n}$ y
            $C=\bigcap_{n=1}^{\infty}C_{n}$ tambien pertenecen a $\mathcal{F}$ por la misma propiedad de cerradura.
            \item[(b)] $B=\left\{\omega\in\Omega:\begin{matrix}
            \omega\in A_{k}\text{ para todos los valores de $k$ salvo, posiblemente}\\
            \text{por un número finito de los mismos}
            \end{matrix}\right\}$\\
            Si $\omega\in B$, entonces existe $n$ tal que $\omega\in B_{n}$. Por definición de $B_{n}$,
            esto implica que $\omega\in A_{k}$ para todo $k\geq n$. Por lo tanto, $\omega$ puede no pertencer
            a lo maximo a los conjuntos $A_{1},\dots,A_{n-1}$, que son finitos. De igual forma, si $\omega$
            pertenece a todos los $A_{k}$ excepto a un numero finito, entonces existe $n$ tal que $\omega\in A_{k}$
            para todo $k\geq n$. Esto implica $\omega\in B_{n}$, y por lo tanto $\omega\in B$.
            \item[(c)] $C=\{\omega\in\Omega:\omega\in A_{k}\text{ para un número infinito de valores de $k$}\}$\\
            Si $\omega\in C$, entonces $\omega\in C_{n}$ para todo $n$. Por definición de $C_{n}$,
            esto significa que para todo $n$ existe $k\geq n$ tal que $\omega\in A_{k}$. Esto
            implica que $\omega$ pertenece a infinitos conjuntos $A_{k}$. De igual manera, si
            $\omega$ pertenece a infinitos $A_{k}$, entonces para todo $n$ existe algun $k\geq n$
            tal que $\omega\in A_{k}$, lo que implica $\omega\in C_{n}$ para todo $n$, y por tanto
            $\omega\in C$.
        \end{enumerate}
    \end{proof}

    \begin{mdframed}[style=mdbluebox,frametitle={Ejercicio 2}]
        Una sucesión de eventos $\{A_k\}_{k=1}^\infty$ en un espacio de probabilidad $(\Omega,\mathcal{F}, P)$ se dice convergente si se cumple que $\liminf_k A_k = \limsup_k A_k$, es decir, su límite inferior y su límite superior coinciden. Demuestre que si $\{A_k\}_{k=1}^\infty$ es convergente con $A = \liminf_k A_k = \limsup_k A_k$, entonces
        \[
        P(A) = \lim_{k \to \infty} P(A_k) .
        \]
    \end{mdframed}

    \begin{proof}
        Por la monotonía de la probabilidad,
        \begin{equation*}
            P(B_{n})\leq P(A_{k})\leq P(C_{n})
        \end{equation*}
        para todo $k\geq n$. Tomando limite inferior y superior cuando $k\to\infty$, se obtiene
        \begin{equation*}
            P(B_{n})\leq\liminf_{k\to\infty}P(A_{k})\leq\limsup_{k\to\infty} P(A_{k})\leq P(C_{n}).
        \end{equation*}
        Ahora, como $\{B_{n}\}$ es creciente,
        \begin{equation*}
            P(A)=P\paren{\bigcup_{n=1}^{\infty}B_{n}}=\lim_{n\to\infty}P(B_{n}).
        \end{equation*}
        Como $\{C_{n}\}$ es decreciente,
        \begin{equation*}
            P(A)=P\paren{\bigcap_{n=1}^{\infty}C_{n}}=\lim_{n\to\infty}P(C_{n}).
        \end{equation*}
        Por lo tanto
        \begin{equation*}
            \lim_{n\to\infty}P(B_{n})=P(A)=\lim_{n\to\infty}P(C_{n}).
        \end{equation*}
        Por ultimo, pasando al limite en la desigualdad
        \begin{equation*}
            P(B_{n})\leq\liminf_{k}P(A_{k})\leq\limsup_{k} P(A_{k})\leq P(C_{n}).
        \end{equation*}
        cuando $n\to\infty$, se obtiene
        \begin{equation*}
            P(A)\leq\liminf_{k}P(A_{k})\leq\limsup_{k} P(A_{k})\leq P(A).
        \end{equation*}
        Esto implica
        \begin{equation*}
            \liminf_{k}P(A_{k})=\limsup_{k}P(A_{k})=P(A),
        \end{equation*}
        y por lo tanto el limite existe y
        \begin{equation*}
            \lim_{n\to\infty}P(A_{k})=P(A).
        \end{equation*}
    \end{proof}

    \begin{mdframed}[style=mdbluebox,frametitle={Ejercicio 3}]
        Sea $(\Omega,\mathcal{F})$ un espacio de medida y suponga que $P_1, P_2$ son dos funciones de probabilidad definidas en este espacio, es decir, $P_1, P_2 : \mathcal{F} \to \mathbb{R}$ son funciones tales que

        \begin{itemize}
            \item[(I)] $P_k(\Omega) = 1$ y $0 \leq P_k(A) \leq 1$ para todo $A \in \mathcal{F}$
            \item[(II)] si $\{A_i\}_{i=1}^\infty \subset \mathcal{F}$ con $A_i \cap A_j = \emptyset$ si $i \neq j$, entonces
            \[
            \sum_{i=1}^{\infty} P(A_i) = P\left(\bigcup_{i=1}^{\infty} A_i\right)
            \]
        \end{itemize}
        con $k = 1, 2$. Demuestre que $\alpha P_1 + (1 - \alpha)P_2$ es también una función de probabilidad. Generalice el resultado para $n$ funciones de probabilidad $P_1, P_2, \dots, P_n$.
    \end{mdframed}

    \begin{mdframed}[style=mdbluebox,frametitle={Ejercicio 4}]
        Para cualesquiera $a, b \in \mathbb{R}$ con $a \leq b$, encuentre los conjuntos
        \[
        \bigcup_{n=1}^{\infty} \left[a + \frac{1}{n}, b - \frac{1}{n}\right]
        \quad \text{y} \quad
        \bigcap_{n=1}^{\infty} \left(a - \frac{1}{n}, b + \frac{1}{n}\right) .
        \]
        Argumente el porqué la $\sigma$-álgebra de $\Omega = [0, 1]$ generada por todos los intervalos cerrados $[a, b] \subset [0, 1]$ coincide con la $\sigma$-álgebra generada por intervalos abiertos $(a, b) \subset [0, 1]$.
    \end{mdframed}
\end{document}